% Options for packages loaded elsewhere
\PassOptionsToPackage{unicode}{hyperref}
\PassOptionsToPackage{hyphens}{url}
\PassOptionsToPackage{dvipsnames,svgnames,x11names}{xcolor}
\documentclass[
  10pt,
  fontset=fandol]{ctexart}
\usepackage{xcolor}
\usepackage[margin=16mm]{geometry}
\usepackage{amsmath,amssymb}
\setcounter{secnumdepth}{-\maxdimen} % remove section numbering
\usepackage{iftex}
\ifPDFTeX
  \usepackage[T1]{fontenc}
  \usepackage[utf8]{inputenc}
  \usepackage{textcomp} % provide euro and other symbols
\else % if luatex or xetex
  \usepackage{unicode-math} % this also loads fontspec
  \defaultfontfeatures{Scale=MatchLowercase}
  \defaultfontfeatures[\rmfamily]{Ligatures=TeX,Scale=1}
\fi
\usepackage{lmodern}
\ifPDFTeX\else
  % xetex/luatex font selection
\fi
% Use upquote if available, for straight quotes in verbatim environments
\IfFileExists{upquote.sty}{\usepackage{upquote}}{}
\IfFileExists{microtype.sty}{% use microtype if available
  \usepackage[]{microtype}
  \UseMicrotypeSet[protrusion]{basicmath} % disable protrusion for tt fonts
}{}
\makeatletter
\@ifundefined{KOMAClassName}{% if non-KOMA class
  \IfFileExists{parskip.sty}{%
    \usepackage{parskip}
  }{% else
    \setlength{\parindent}{0pt}
    \setlength{\parskip}{6pt plus 2pt minus 1pt}}
}{% if KOMA class
  \KOMAoptions{parskip=half}}
\makeatother
\usepackage{longtable,booktabs,array}
\newcounter{none} % for unnumbered tables
\usepackage{calc} % for calculating minipage widths
% Correct order of tables after \paragraph or \subparagraph
\usepackage{etoolbox}
\makeatletter
\patchcmd\longtable{\par}{\if@noskipsec\mbox{}\fi\par}{}{}
\makeatother
% Allow footnotes in longtable head/foot
\IfFileExists{footnotehyper.sty}{\usepackage{footnotehyper}}{\usepackage{footnote}}
\makesavenoteenv{longtable}
\setlength{\emergencystretch}{3em} % prevent overfull lines
\providecommand{\tightlist}{%
  \setlength{\itemsep}{0pt}\setlength{\parskip}{0pt}}
\usepackage{amsmath,amssymb,mathtools,needspace}
\xeCJKDeclareCharClass{CJK}{"2032,"0400->"04FF,"2160->"216B,"2460->"2473,"25A0,"25C7,"25CB,"25CF}
\IfFontExistsTF{Arial Unicode MS}{\xeCJKsetup{AutoFallBack=true}\setCJKfallbackfamilyfont{\CJKrmdefault}{Arial Unicode MS}}{}
\setcounter{secnumdepth}{0}
\setlength{\emergencystretch}{2em}
\allowdisplaybreaks
\usepackage{bookmark}
\IfFileExists{xurl.sty}{\usepackage{xurl}}{} % add URL line breaks if available
\urlstyle{same}
\hypersetup{
  pdftitle={计算步骤与例题},
  colorlinks=true,
  linkcolor={Maroon},
  filecolor={Maroon},
  citecolor={Blue},
  urlcolor={Blue},
  pdfcreator={LaTeX via pandoc}}

\title{计算步骤与例题}
\author{}
\date{}

\begin{document}
\maketitle

\subsubsection{2.8.4 计算步骤与例题}\label{ux8ba1ux7b97ux6b65ux9aa4ux4e0eux4f8bux9898}

样条函数，特别是三次样条在实际中有广泛的应用，在计算机上也容易实现。下面以方程（2.8.12）为例，说明在计算机上求 \(S(x)\) 的\textbf{算法步骤}：

\textbf{步 1}　输入初始数据 \(x_j,y_j\)（\(j=0,1,\cdots,n\)）及 \(f'_0,f'_n\) 和 \(n\)。

\textbf{步 2}　\(j\) 从 0 到 \(n-1\) 计算 \(h_j=x_{j+1}-x_j\) 及 \(f[x_j,x_{j+1}]\)。

\textbf{步 3}　\(j\) 从 1 到 \(n-1\) 由式（2.8.10）及式（2.8.11）计算 \(\lambda_j,\mu_j,g_j\)。

\textbf{步 4}　用追赶法（公式见 7.4.3 节）解方程（2.8.12），求出 \(m_j\)（\(j=1,2,\cdots,n-1\)）。

\textbf{步 5}　计算 \(S(x)\) 的系数或计算 \(S(x)\) 在若干点上的值，并打印结果。

\textbf{步 6}　给定函数 \(f(x)=\frac1{1+x^2},-5\leqslant x\leqslant5\)，节点 \(x_k=-5+k\)（\(k=0,1,\cdots,10\)），用三次样条插值求 \(S_{10}(x)\)。取 \(S_{10}(x_k)=f(x_k)\)（\(k=0,1,\cdots,10\)），有

\[S'_{10}(-5)=f'(-5),\quad S'_{10}(5)=f'(5).\]

利用上述步骤编制的程序计算 \(S_{10}(x)\) 在表 2.8 所列各点的值，并与 \(f(x)\) 及 2.7 节计算出的 \(L_{10}(x)\) 比较，结果也如表 2.8 所示。由表 2.8 可看出，用三次样条插值得到的 \(S_{10}(x)\) 能很好地逼近 \(f(x)\)，不会出现 Lagrange 插值 \(L_{10}(x)\) 的``Runge''现象。

\href{../../source-images/pdf-054.jpeg}{查看原页}

\begin{quote}
本页接上页 2.8.4 正文。
\end{quote}

\(f(x)\) 与 \(L_{10}(x)\) 的图形如图 2.5 所示，如画出 \(S_{10}(x)\) 的曲线，将与 \(f(x)\) 很近似。

\textbf{表 2.8}

{\def\LTcaptype{none} % do not increment counter
\begin{longtable}[]{@{}
  >{\raggedright\arraybackslash}p{(\linewidth - 14\tabcolsep) * \real{0.1250}}
  >{\raggedright\arraybackslash}p{(\linewidth - 14\tabcolsep) * \real{0.1250}}
  >{\raggedright\arraybackslash}p{(\linewidth - 14\tabcolsep) * \real{0.1250}}
  >{\raggedright\arraybackslash}p{(\linewidth - 14\tabcolsep) * \real{0.1250}}
  >{\raggedright\arraybackslash}p{(\linewidth - 14\tabcolsep) * \real{0.1250}}
  >{\raggedright\arraybackslash}p{(\linewidth - 14\tabcolsep) * \real{0.1250}}
  >{\raggedright\arraybackslash}p{(\linewidth - 14\tabcolsep) * \real{0.1250}}
  >{\raggedright\arraybackslash}p{(\linewidth - 14\tabcolsep) * \real{0.1250}}@{}}
\toprule\noalign{}
\begin{minipage}[b]{\linewidth}\raggedright
\(x\)
\end{minipage} & \begin{minipage}[b]{\linewidth}\raggedright
\(\frac1{1+x^2}\)
\end{minipage} & \begin{minipage}[b]{\linewidth}\raggedright
\(S_{10}(x)\)
\end{minipage} & \begin{minipage}[b]{\linewidth}\raggedright
\(L_{10}(x)\)
\end{minipage} & \begin{minipage}[b]{\linewidth}\raggedright
\(x\)
\end{minipage} & \begin{minipage}[b]{\linewidth}\raggedright
\(\frac1{1+x^2}\)
\end{minipage} & \begin{minipage}[b]{\linewidth}\raggedright
\(S_{10}(x)\)
\end{minipage} & \begin{minipage}[b]{\linewidth}\raggedright
\(L_{10}(x)\)
\end{minipage} \\
\midrule\noalign{}
\endhead
\bottomrule\noalign{}
\endlastfoot
−5.0 & 0.038 46 & 0.038 46 & 0.038 46 & −2.3 & 0.158 98 & 0.161 15 & 0.241 45 \\
−4.8 & 0.041 60 & 0.037 58 & 1.804 38 & −2.0 & 0.200 00 & 0.200 00 & 0.200 00 \\
−4.5 & 0.047 06 & 0.042 48 & 1.578 72 & −1.8 & 0.235 85 & 0.231 54 & 0.188 78 \\
−4.3 & 0.051 31 & 0.048 42 & 0.888 08 & −1.5 & 0.307 69 & 0.297 44 & 0.235 35 \\
−4.0 & 0.058 82 & 0.058 82 & 0.058 82 & −1.3 & 0.371 75 & 0.361 33 & 0.316 50 \\
−3.8 & 0.064 77 & 0.065 56 & −0.201 30 & −1.0 & 0.500 00 & 0.500 00 & 0.500 00 \\
−3.5 & 0.075 47 & 0.076 06 & −0.226 20 & −0.8 & 0.609 76 & 0.624 20 & 0.643 16 \\
−3.3 & 0.084 10 & 0.084 26 & −0.108 32 & −0.5 & 0.800 00 & 0.820 51 & 0.843 40 \\
−3.0 & 0.100 00 & 0.100 00 & 0.100 00 & −0.3 & 0.917 43 & 0.927 54 & 0.940 90 \\
−2.8 & 0.113 12 & 0.113 66 & 0.198 37 & 0 & 1.000 00 & 1.000 00 & 1.000 00 \\
−2.5 & 0.137 93 & 0.139 71 & 0.253 76 & & & & \\
\end{longtable}
}

\begin{quote}
表格数值及原有小数分组空格均按扫描转录；\href{../assets/ch02b/table-2-8.png}{表 2.8 原图裁切}。
\end{quote}

\begin{center}\rule{0.5\linewidth}{0.5pt}\end{center}

\subsection{相关原页校注（agent 补充）}\label{ux76f8ux5173ux539fux9875ux6821ux6ce8agent-ux8865ux5145}

以下校注来自本节覆盖的原页，可能也涉及同页相邻小节；原文照录与校正说明分开保留。

\subsubsection{原书 PDF 页 54 的校注}\label{ux539fux4e66-pdf-ux9875-54-ux7684ux6821ux6ce8}

\href{../pages/pdf-054.md}{完整原页转录} · \href{../../source-images/pdf-054.jpeg}{原页图像}

校注记录：ch02b-E006, ch02b-E007, ch02b-E008。

\begin{quote}
\textbf{校注（agent 补充；原书疑误）}：式（2.8.21）原图在 \(\max\) 下用 \(i\)，而被取最大值的分量印为 \(x_j\)；已原样保留，标准表达需将哑指标统一。式（2.8.23）括号内第一项原图印为 \(|a_{ij}|\)；结合（2.8.22）的对角占优量，应为 \(|a_{ii}|\)，并对行指标 \(i\) 取最小值。参见\href{../assets/ch02b/pdf-054-norm-detail.png}{范数公式原图裁切}。
\end{quote}

\begin{quote}
\textbf{校注（agent 补充；原书表值与条件不一致）}：表 2.8 的 \(S_{10}\) 数值在此按原图保留。按上一页明定的节点、\(f(x)=1/(1+x^2)\) 及 \(S'_{10}(\pm5)=f'(\pm5)\)，精确有理数解三转角方程给出 \(S_{10}(-4.8)\approx0.04162183\)、\(S_{10}(-4.5)\approx0.04716801\)，与原表的 \(0.03758\)、\(0.04248\) 不同；这不是 OCR 差异。复核结果见\href{../../staging/ch02b/table-2-8-check.txt}{独立数值检查}。本校注不以复算值替换原表。
\end{quote}

\subsection{本节覆盖原页的校注记录}\label{ux672cux8282ux8986ux76d6ux539fux9875ux7684ux6821ux6ce8ux8bb0ux5f55}

下列记录按原页关联，含同页相邻小节的校注；计算时同时读取上文校注和完整原页。

\begin{itemize}
\tightlist
\item
  \texttt{ch02b-E006}：原书 PDF 页 54
\item
  \texttt{ch02b-E007}：原书 PDF 页 54
\item
  \texttt{ch02b-E008}：原书 PDF 页 54
\end{itemize}

\end{document}
